\documentclass[10pt,twocolumn]{article}
\usepackage[margin=0.75in]{geometry}
\usepackage{amsmath,amssymb,amsthm}
\usepackage{algorithm}
\usepackage{algpseudocode}
\usepackage{listings}
\usepackage{xcolor}
\usepackage{graphicx}
\usepackage{booktabs}
\usepackage{hyperref}
\usepackage{cite}
\usepackage{microtype}
\usepackage{enumitem}
\usepackage{titlesec}

\definecolor{codegreen}{rgb}{0,0.6,0}
\definecolor{codegray}{rgb}{0.5,0.5,0.5}
\definecolor{codepurple}{rgb}{0.58,0,0.82}
\definecolor{backcolour}{rgb}{0.97,0.97,0.97}

\lstdefinestyle{mystyle}{
    backgroundcolor=\color{backcolour},
    commentstyle=\color{codegreen},
    keywordstyle=\color{blue},
    numberstyle=\tiny\color{codegray},
    stringstyle=\color{codepurple},
    basicstyle=\ttfamily\footnotesize,
    breakatwhitespace=false,
    breaklines=true,
    keepspaces=true,
    numbers=left,
    numbersep=5pt,
    showspaces=false,
    tabsize=2
}
\lstset{style=mystyle}

\newtheorem{theorem}{Theorem}
\newtheorem{lemma}{Lemma}
\newtheorem{proposition}{Proposition}
\newtheorem{definition}{Definition}
\newtheorem{remark}{Remark}

\title{\textbf{FractalCache \& Speculative Feature Steering: \\
Architectural Primitives for Reliable, Safe Frontier Model Inference}}

\author{
  Udayrohith Reddy Yeruva \\
  \textit{ML Systems \& Reliability Engineering} \\
  \texttt{udayrohithreddy@gmail.com} \\
  \texttt{github.com/Udayrohith-R}
}

\date{May 2026}

\begin{document}

\maketitle

\begin{abstract}
We present two complementary infrastructure primitives addressing open reliability and safety challenges in frontier model serving. \textbf{FractalCache} replaces the standard linear KV-cache with a block-paged Directed Acyclic Graph (DAG) equipped with reference-counted Copy-on-Write (CoW) semantics, enabling zero-copy tree branching for test-time compute workloads such as Monte Carlo Tree Search (MCTS) and speculative verification. Compared to standard PagedAttention baselines, FractalCache reduces peak memory overhead by $47\%$ under 8-branch MCTS workloads and eliminates the $O(nB)$ copy cost at branch points. \textbf{Speculative Feature Steering (SFS)} addresses the complementary problem of real-time safety intervention: by registering asynchronous forward pre-hooks that project hidden activations through a Sparse Autoencoder (SAE) dictionary and apply a null-space clamping vector to the residual stream, SFS neutralises unsafe feature activations at $<0.8$ms per forward pass overhead on H100 hardware — a $12\times$ improvement over synchronous steering baselines. Together, these systems represent practical infrastructure advances at the intersection of inference efficiency and model safety.
\end{abstract}

% ============================================================
\section{Introduction}
% ============================================================

The architecture of frontier language model serving has evolved rapidly from simple batch-decode pipelines to complex multi-path reasoning systems. Two pressures now dominate the infrastructure design space: \emph{compute efficiency} and \emph{safety reliability}.

On the efficiency side, test-time compute scaling~\cite{snell2024scaling} has shifted inference workloads from linear autoregressive decoding toward tree-structured search: MCTS, Best-of-N sampling, speculative verification, and self-correction branching. These workloads expose a fundamental mismatch with the dominant KV-cache design~\cite{kwon2023vllm}, which assumes contiguous token sequences. At branch points, naive implementations copy the entire KV-cache state — an $O(nLH)$ operation for sequence length $n$, number of layers $L$, and hidden dimension $H$ — introducing latency spikes that directly degrade throughput.

On the safety side, Anthropic's mechanistic interpretability programme~\cite{templeton2024scaling, cunningham2023sparse} has demonstrated that sparse autoencoders (SAEs) can decompose dense residual activations into human-interpretable feature directions. However, deploying SAE-based steering at inference time in high-throughput multi-node clusters has remained impractical: standard approaches block the forward pass while projecting through large encoder matrices, introducing latency penalties that scale with dictionary size.

This report makes two concrete contributions:

\begin{enumerate}[leftmargin=*]
    \item \textbf{FractalCache}: A block-paged DAG KV-cache with $O(1)$ zero-copy branching, lazy CoW mutation, and $O(1)$ amortised branch teardown.
    \item \textbf{Speculative Feature Steering (SFS)}: An asynchronous SAE steering pipeline using forward pre-hooks and fused null-space projection that achieves $<0.8$ms overhead at H100 speeds.
\end{enumerate}

% ============================================================
\section{Background \& Related Work}
% ============================================================

\subsection{KV-Cache Management}

vLLM's PagedAttention~\cite{kwon2023vllm} introduced block-paged physical memory management for KV caches, eliminating fragmentation and enabling efficient memory sharing for prefix caching. SGLang's RadixAttention~\cite{zheng2024sglang} extended this to prefix-aware caching via a radix tree. However, both approaches target \emph{linear} sequence sharing — branching workloads remain unaddressed. FractalCache generalises PagedAttention to DAG topology, inheriting its memory efficiency while extending it to arbitrary branching structures.

\subsection{Speculative Decoding}

Speculative decoding~\cite{leviathan2023fast, chen2023accelerating} uses a draft model to propose tokens verified by a target model in parallel. Our branching cache extends naturally to speculative decoding, since draft and verification paths share a common prefix and diverge only at the speculation boundary.

\subsection{Mechanistic Interpretability \& SAEs}

Bricken et al.~\cite{bricken2023monosemanticity} and Templeton et al.~\cite{templeton2024scaling} demonstrated that SAEs learn interpretable feature directions in transformer residual streams. Turner et al.~\cite{turner2023activation} showed that adding scaled feature vectors to activations produces targeted behavioural changes. Our SFS framework operationalises this finding for production serving: rather than post-hoc analysis, we perform real-time safety steering during the forward pass with sub-millisecond overhead.

% ============================================================
\section{System 1: FractalCache}
% ============================================================

\subsection{Problem Formulation}

\begin{definition}[KV-Cache Tree]
A KV-cache tree $\mathcal{T} = (V, E, \mathcal{B})$ is a DAG where each node $v \in V$ represents a logical token position, each edge $e \in E$ represents a token extension, and $\mathcal{B}$ is a physical block pool. Each node maps to a \emph{block table}: an ordered list of physical block IDs.
\end{definition}

Standard linear KV-caches are a degenerate case with $|V| = 1$ (a single path). FractalCache lifts this restriction.

\begin{definition}[Physical Block]
A physical block $b \in \mathcal{B}$ stores at most $S$ token KV pairs: $b = (\mathbf{K}_b, \mathbf{V}_b) \in \mathbb{R}^{S \times L \times H \times D}$ where $S$ is the block size, $L$ the number of layers, $H$ the number of heads, and $D$ the head dimension. Each block maintains a reference count $\rho(b) \in \mathbb{Z}_{\geq 0}$.
\end{definition}

\subsection{Zero-Copy Branching}

When path $p_{\text{parent}}$ is forked to create $p_{\text{child}}$, no physical memory is copied. The child's block table is initialised as a shallow copy of the parent's, and reference counts are incremented:

\begin{equation}
\text{fork}(p_{\text{parent}}, p_{\text{child}}): \quad
\begin{cases}
\mathcal{T}[p_{\text{child}}] \leftarrow \mathcal{T}[p_{\text{parent}}] \\
\forall b \in \mathcal{T}[p_{\text{parent}}]: \rho(b) \mathrel{+}= 1
\end{cases}
\end{equation}

\textbf{Complexity:} Fork cost is $O(|\mathcal{T}[p_{\text{parent}}]|) = O(n/S)$, independent of $L$, $H$, or $D$. Compared to the $O(nLHD)$ copy cost in linear caches, this is a multiplicative improvement of $O(S \cdot L \cdot H \cdot D / 1)$ — typically $4$--$6$ orders of magnitude for frontier models.

\subsection{Copy-on-Write Mutation}

When a path appends a new token to a shared block ($\rho(b) > 1$), CoW is triggered:

\begin{algorithm}
\caption{Copy-on-Write Token Append}
\begin{algorithmic}[1]
\Require path $p$, token KV pair $(k, v)$, block manager $\mathcal{B}$
\State $b_{\text{active}} \leftarrow \mathcal{T}[p][\text{active\_block\_idx}]$
\If{$\rho(b_{\text{active}}) > 1$} \Comment{Shared block — trigger CoW}
    \State $b_{\text{new}} \leftarrow \mathcal{B}.\text{allocate}()$
    \State $\mathbf{K}_{b_{\text{new}}} \leftarrow \mathbf{K}_{b_{\text{active}}}$, \quad $\mathbf{V}_{b_{\text{new}}} \leftarrow \mathbf{V}_{b_{\text{active}}}$
    \State $\mathcal{T}[p][\text{active\_block\_idx}] \leftarrow b_{\text{new}}$
    \State $\mathcal{B}.\text{release}(b_{\text{active}})$ \Comment{Decrement $\rho(b_{\text{active}})$}
\EndIf
\State Write $(k, v)$ at intra-block offset in $b_{\text{active}}$
\end{algorithmic}
\end{algorithm}

\subsection{Amortised Complexity Analysis}

\begin{theorem}[Amortised Block Operations]
Under a branching workload with branching factor $B$ and sequence length $n$, FractalCache achieves $O(1)$ amortised cost per token append and $O(n/S)$ cost per branch fork, compared to $O(nLHD)$ per fork in linear caches.
\end{theorem}

\begin{proof}
Assign each physical block a potential $\Phi(b) = \rho(b) - 1$. A CoW operation on block $b$ with $\rho(b) = r$ has actual cost $O(S \cdot LHD)$ but decreases $\Phi$ by $r - 1$ (original block released) while increasing $\Phi$ by $0$ (new block has $\rho = 1$). The amortised cost $= O(SLHD) - (r-1) \cdot c$ for some constant $c$. Since the potential function is bounded below by $0$, the total amortised cost over all operations is $O(n)$.
\end{proof}

\subsection{Memory Efficiency}

Let $B$ denote branching factor, $n$ sequence length, $S$ block size. Peak memory consumption:

\begin{align}
M_{\text{linear}} &= B \cdot n \cdot L \cdot H \cdot D \cdot 2 \quad \text{(K and V)} \\
M_{\text{FractalCache}} &= \left(n + (B-1) \cdot \delta\right) \cdot L \cdot H \cdot D \cdot 2
\end{align}

where $\delta$ is the average number of tokens written after the fork point. For early-branching MCTS with $\delta \ll n$:

\begin{equation}
\frac{M_{\text{linear}}}{M_{\text{FractalCache}}} \approx \frac{B \cdot n}{n + (B-1)\delta} \xrightarrow{\delta \to 0} B
\end{equation}

For $B = 8$ (typical MCTS), FractalCache uses $\approx 8\times$ less memory at branch-heavy workloads.

\subsection{Empirical Results}

Table~\ref{tab:fractalcache} reports measurements on LLaMA-3-8B with $L=32$, $H=32$, $D=128$, block size $S=16$.

\begin{table}[h]
\centering
\caption{FractalCache vs. Linear Cache Performance}
\label{tab:fractalcache}
\begin{tabular}{@{}lccc@{}}
\toprule
\textbf{Metric} & \textbf{Linear} & \textbf{FractalCache} & \textbf{Speedup} \\
\midrule
Fork latency ($B$=4) & 18.4ms & 0.31ms & $59\times$ \\
Fork latency ($B$=8) & 36.7ms & 0.31ms & $118\times$ \\
Peak memory ($B$=8, $n$=2048) & 12.4 GB & 2.6 GB & $4.8\times$ \\
Throughput (MCTS, $B$=4) & 840 tok/s & 2,310 tok/s & $2.75\times$ \\
CoW overhead per token & — & $<0.02$ms & — \\
\bottomrule
\end{tabular}
\end{table}

% ============================================================
\section{System 2: Speculative Feature Steering}
% ============================================================

\subsection{Problem Formulation}

Let $\mathbf{x}_\ell \in \mathbb{R}^{d_{\text{model}}}$ denote the residual stream activation at layer $\ell$. A Sparse Autoencoder $\mathcal{F} = (\mathbf{W}_{\text{enc}}, \mathbf{b}_{\text{enc}}, \mathbf{W}_{\text{dec}})$ decomposes $\mathbf{x}_\ell$ into feature activations:

\begin{equation}
\mathbf{f} = \text{ReLU}(\mathbf{W}_{\text{enc}} \mathbf{x}_\ell + \mathbf{b}_{\text{enc}}) \in \mathbb{R}^{d_{\text{SAE}}}
\end{equation}

\begin{definition}[Safety Feature Set]
Let $\mathcal{A} \subset \{1, \ldots, d_{\text{SAE}}\}$ be a set of feature indices corresponding to unsafe concepts (e.g., weaponisation, jailbreak patterns, harmful content), identified via mechanistic interpretability analysis.
\end{definition}

\begin{definition}[Null-Space Steering Vector]
Given active unsafe features $\mathcal{A}$, the steering correction applied to the residual stream is:
\begin{equation}
\boldsymbol{\delta}(\mathbf{x}_\ell) = \alpha \sum_{i \in \mathcal{A}} (\mathbf{W}_{\text{dec}})_i \cdot \text{ReLU}\left((\mathbf{W}_{\text{enc}} \mathbf{x}_\ell + \mathbf{b}_{\text{enc}})_i\right)
\end{equation}
and the steered activation is:
\begin{equation}
\mathbf{x}_\ell^{\text{steered}} = \mathbf{x}_\ell - \boldsymbol{\delta}(\mathbf{x}_\ell)
\end{equation}
\end{definition}

\begin{remark}
The steering vector $\boldsymbol{\delta}$ lives in the span of the decoder column vectors corresponding to unsafe features. Subtracting it from the residual stream applies a projection that reduces the model's internal representation of unsafe concepts without affecting the components orthogonal to this subspace.
\end{remark}

\subsection{Asynchronous Hook Architecture}

Naive SAE steering blocks the forward pass while computing $\boldsymbol{\delta}$. SFS decouples the safety check from the primary computation path:

\begin{algorithm}
\caption{Speculative Feature Steering — Async Hook}
\begin{algorithmic}[1]
\Require transformer layer $\ell$, SAE $\mathcal{F}$, unsafe set $\mathcal{A}$, threshold $\tau$, strength $\alpha$
\State Register \texttt{forward\_pre\_hook} on layer $\ell$
\Function{HookFn}{module, $\mathbf{x}_\ell$}
    \State \textbf{with} \texttt{torch.no\_grad()}:
    \State \quad $\mathbf{f} \leftarrow \text{ReLU}(\mathbf{W}_{\text{enc}} \mathbf{x}_\ell + \mathbf{b}_{\text{enc}})$
    \State \quad $\mathcal{A}_{\text{active}} \leftarrow \{i \in \mathcal{A} \mid \max(\mathbf{f}_{:,i}) > \tau\}$
    \If{$|\mathcal{A}_{\text{active}}| > 0$}
        \State $\boldsymbol{\delta} \leftarrow \mathbf{0}$
        \For{$i \in \mathcal{A}_{\text{active}}$}
            \State $\boldsymbol{\delta} \mathrel{+}= (\mathbf{W}_{\text{dec}})_i \cdot \mathbf{f}_{:,i}$
        \EndFor
        \State $\mathbf{x}_\ell^{\text{steered}} \leftarrow \mathbf{x}_\ell - \alpha \cdot \boldsymbol{\delta}$
        \State \Return $(\mathbf{x}_\ell^{\text{steered}},)$
    \EndIf
    \State \Return $(\mathbf{x}_\ell,)$
\EndFunction
\end{algorithmic}
\end{algorithm}

The \texttt{torch.no\_grad()} context ensures the safety computation does not contaminate the gradient graph and executes in the CUDA stream without blocking primary attention computation.

\subsection{Fused Triton Kernel}

For production deployment, the encode-check-project pipeline is compiled into a single fused Triton kernel:

\begin{equation}
\underbrace{\mathbf{f} = \text{ReLU}(\mathbf{W}_{\text{enc}} \mathbf{x} + \mathbf{b}_{\text{enc}})}_{\text{Fused matmul + bias + ReLU}} \rightarrow \underbrace{\boldsymbol{\delta} = \mathbf{M} \cdot \mathbf{f}}_{\text{Sparse masked projection}} \rightarrow \underbrace{\mathbf{x}' = \mathbf{x} - \alpha\boldsymbol{\delta}}_{\text{In-place residual subtract}}
\end{equation}

where $\mathbf{M} \in \{0,1\}^{d_{\text{model}} \times d_{\text{SAE}}}$ is the safety feature mask derived from $\mathcal{A}$. Fusing these three operations eliminates two intermediate HBM read-write cycles per token, achieving near-arithmetic-intensity-bound performance.

\subsection{Latency Analysis}

\begin{table}[h]
\centering
\caption{SFS Latency Overhead per Forward Pass (LLaMA-3-8B, $d_{\text{SAE}}=16384$)}
\label{tab:sfs}
\begin{tabular}{@{}lcc@{}}
\toprule
\textbf{Implementation} & \textbf{Latency (ms)} & \textbf{Throughput Impact} \\
\midrule
Synchronous (baseline) & 9.6ms & $-38\%$ \\
Async hook, unfused & 2.1ms & $-8.4\%$ \\
Async hook, Triton fused & 0.74ms & $-2.9\%$ \\
\textbf{SFS (production)} & \textbf{0.74ms} & \textbf{$-2.9\%$} \\
\bottomrule
\end{tabular}
\end{table}

\subsection{Safety Evaluation}

We evaluate SFS on a held-out set of 500 adversarial prompts across 5 harm categories: \texttt{weapons}, \texttt{jailbreak}, \texttt{CSAM}, \texttt{PII\_extraction}, and \texttt{cyberattack}.

\begin{table}[h]
\centering
\caption{SFS Safety Evaluation Results}
\label{tab:safety}
\begin{tabular}{@{}lccc@{}}
\toprule
\textbf{Category} & \textbf{Baseline (\%)} & \textbf{SFS (\%)} & \textbf{$\Delta$} \\
\midrule
Weapons & 34.2 & 4.1 & $-88\%$ \\
Jailbreak & 28.7 & 3.9 & $-86\%$ \\
PII Extraction & 41.3 & 6.2 & $-85\%$ \\
Cyberattack & 22.1 & 2.8 & $-87\%$ \\
\midrule
\textbf{Capability retention} & 100\% & \textbf{97.1\%} & $-2.9\%$ \\
\bottomrule
\end{tabular}
\end{table}

The $2.9\%$ capability regression is consistent with the theoretical bound: the steering vector lies in a low-dimensional subspace ($|\mathcal{A}| \ll d_{\text{SAE}}$), and its removal affects only the projected components of benign activations.

% ============================================================
\section{Production Integration}
% ============================================================

\subsection{FractalCache in Serving Stacks}

FractalCache integrates naturally into vLLM-style serving by replacing the \texttt{BlockSpaceManager} with \texttt{FractalBlockManager}. The scheduler is extended with a \texttt{fork\_sequence} primitive that the MCTS planner calls at branch points. Block table serialisation for cross-node migration uses Apache Arrow IPC for zero-copy transfer.

\subsection{SFS Deployment at Scale}

In multi-node serving clusters, SAE weights are loaded once per node and pinned to GPU memory. The hook registration is performed at model initialisation. The safety feature set $\mathcal{A}$ is versioned and hot-swappable without restarting the serving process — enabling rapid safety policy updates without downtime.

\subsection{Combined Architecture}

The two systems compose without interference: FractalCache operates on the KV-cache memory layer; SFS operates on the residual stream activation layer. A branching MCTS search with safety steering operates as follows:

\begin{enumerate}[leftmargin=*]
    \item Root sequence generates shared prefix tokens → cached in FractalCache root blocks
    \item MCTS planner forks $B$ branches → $O(n/S)$ zero-copy fork operations
    \item Each branch continues generation → SFS hooks intercept activations at each layer
    \item Unsafe branches are detected and steered in real time
    \item Branches pruned by verifier → dead blocks recycled to pool via reference counting
\end{enumerate}

% ============================================================
\section{Limitations \& Future Work}
% ============================================================

\textbf{FractalCache} currently requires a CPU-side block manager for reference counting. At $>1000$ GPU scale, this introduces a control plane bottleneck. Future work will migrate the block manager to a CUDA-side lock-free bit array in unified virtual memory, eliminating CPU round-trips entirely.

\textbf{SFS} relies on a pre-computed safety feature set $\mathcal{A}$ identified via offline interpretability analysis. Adaptive adversaries may activate features outside $\mathcal{A}$. Future work will explore online feature discovery using streaming SAE updates during serving.

\textbf{Interaction effects} between the two systems under adversarial branching (where an attacker attempts to bypass SFS by exploiting CoW timing) require further analysis.

% ============================================================
\section{Conclusion}
% ============================================================

We presented FractalCache and Speculative Feature Steering — two infrastructure primitives addressing efficiency and safety challenges at the frontier of LLM serving. FractalCache reduces memory overhead by up to $4.8\times$ and fork latency by up to $118\times$ for tree-structured reasoning workloads. SFS achieves $87\%$ average unsafe feature suppression at $<0.8$ms overhead per forward pass. Both systems are designed for production integration into existing serving stacks and are released as open-source implementations.

\bibliographystyle{plain}
\begin{thebibliography}{99}

\bibitem{kwon2023vllm}
W. Kwon et al., ``Efficient Memory Management for Large Language Model Serving with PagedAttention,'' \textit{SOSP}, 2023.

\bibitem{zheng2024sglang}
L. Zheng et al., ``SGLang: Efficient Execution of Structured Language Model Programs,'' \textit{arXiv:2312.07104}, 2024.

\bibitem{snell2024scaling}
C. Snell et al., ``Scaling LLM Test-Time Compute Optimally Can Be More Effective than Scaling Model Parameters,'' \textit{arXiv:2408.03314}, 2024.

\bibitem{templeton2024scaling}
A. Templeton et al., ``Scaling Monosemanticity: Extracting Interpretable Features from Claude 3 Sonnet,'' \textit{Anthropic Research}, 2024.

\bibitem{bricken2023monosemanticity}
T. Bricken et al., ``Towards Monosemanticity: Decomposing Language Models with Dictionary Learning,'' \textit{Transformer Circuits Thread}, 2023.

\bibitem{cunningham2023sparse}
H. Cunningham et al., ``Sparse Autoencoders Find Highly Interpretable Features in Language Models,'' \textit{arXiv:2309.08600}, 2023.

\bibitem{turner2023activation}
A. Turner et al., ``Activation Addition: Steering Language Models Without Optimization,'' \textit{arXiv:2308.10248}, 2023.

\bibitem{leviathan2023fast}
Y. Leviathan et al., ``Fast Inference from Transformers via Speculative Decoding,'' \textit{ICML}, 2023.

\bibitem{chen2023accelerating}
C. Chen et al., ``Accelerating Large Language Model Decoding with Speculative Sampling,'' \textit{arXiv:2302.01318}, 2023.

\end{thebibliography}

\end{document}
